\chapter{The shared framework}
\label{chap:framework}

Deligne's conjecture is a single assertion instantiated at many different classes of
motive. This chapter isolates the data and the assertions that every instance shares,
so that GL(1), GL(2), Hecke characters and the rest are \emph{instances} rather than
restatements. Nothing here is specific to any one case.

Mathlib has no category of motives, and we do not build one. The structure below
abstracts exactly the six pieces of data an instance must supply, with no motivic
content whatsoever, and characterises the expected behaviour of that data through four
requirements.

\section{The data}

\begin{definition}[Deligne package]
  \label{def:pkg}
  \lean{Deligne.Pkg}
  \leanok
  Let $\Omega$ be a type, thought of as indexing a family of pure motives over $\Q$.
  A \emph{Deligne package} on $\Omega$ is a tuple $(L, E, w, \gamma, c^{+}, \gal)$ consisting of
  \begin{enumerate}
    \item $L \colon \Omega \to \C \to \C$, the analytic $L$-function $L(M,s)$;
    \item $E \colon \Omega \to \{\text{intermediate fields } \Q \subseteq E \subseteq \C\}$,
      the value field $E(M)$;
    \item $w \in \Z$, the common weight: the functional equation exchanges $s$ and $w+1-s$;
    \item $\mathrm{shifts} \colon \Omega \to \operatorname{Multiset}(\Z)$, the shifts of the
      archimedean Gamma factor, which is then
      $\gamma(M,s) \coloneqq \prod_{a \in \mathrm{shifts}(M)} \Gamma_{\R}(s+a)$;
    \item $c^{+} \colon \Omega \to \Z \to \C$, Deligne's period $c^{+}(M(n))$;
    \item $\gal \colon \operatorname{Gal}(\C/\Q) \longrightarrow \operatorname{Sym}(\Omega)$,
      a group \emph{homomorphism} into the permutations of $\Omega$, giving the action
      $\sigma \mapsto (M \mapsto M^{\sigma})$. Here $\operatorname{Gal}(\C/\Q)$ is the group
      of $\Q$-algebra automorphisms of $\C$, which is every field automorphism of $\C$,
  \end{enumerate}
  subject to four requirements:
  \begin{enumerate}
    \setcounter{enumi}{6}
    \item conjugate objects have the same Hodge numbers:
      $\mathrm{shifts}(M^{\sigma}) = \mathrm{shifts}(M)$;
    \item the value field transports along the action:
      $E(M^{\sigma}) = \sigma\!\left(E(M)\right)$;
    \item the value field is algebraic: every $x \in E(M)$ is algebraic over $\Q$;
    \item periods do not vanish at critical integers:
      $\forall M,\ \forall n,\ \operatorname{ord}_{n}\gamma(M,\cdot) \geq 0 \wedge
      \operatorname{ord}_{w+1-n}\gamma(M,\cdot) \geq 0 \implies c^{+}(M(n)) \neq 0$.
  \end{enumerate}
  In the literature the objects of $\Omega$ are motives and most of these fields are derived
  from them; in lieu of a formal definition of a motive, the requirements record the behaviour
  the derived data would have. In Lean, (x) is stated with criticality written out as arithmetic
  on the shifts (Lemma~\ref{lem:prodGammaR}), because Definition~\ref{def:isCritical} is made
  from a package and so cannot be named inside the package's own definition.
\end{definition}

\begin{definition}[Critical integers]
  \label{def:isCritical}
  \lean{Deligne.Pkg.IsCritical, Deligne.Pkg.gammaFactor, Deligne.Pkg.isCritical\_iff}
  \leanok
  \uses{def:pkg, lem:prodGammaR}
  $n$ is \emph{critical} for $M$ when $\gamma(M,\cdot)$ has a pole at neither $n$ nor $w+1-n$ ---
  the two points exchanged by the functional equation. ``No pole at $s$'' is
  $\operatorname{ord}_{s}\gamma(M,\cdot) \geq 0$, the meromorphic order in mathlib's sense
  (\lean{meromorphicOrderAt}). By Lemma~\ref{lem:prodGammaR} this is the condition that $n + a$
  and $w + 1 - n + a$ are odd or positive for every shift $a$, which is the form an instance
  checks it in and the form requirement (x) of Definition~\ref{def:pkg} is stated in.
\end{definition}

\begin{lemma}[$\Gamma_{\R}$ has no zeros]
  \label{lem:GammaR-pole}
  \lean{Complex.meromorphicOrderAt\_Gammaℝ\_add\_nonpos,
    Complex.meromorphicOrderAt\_Gammaℝ\_add\_intCast\_nonneg\_iff}
  \leanok
  For $a, s \in \C$ the order of $z \mapsto \Gamma_{\R}(z+a)$ at $s$ is at most $0$: every point
  is either a regular point or a pole. At an integer $m$, with $a \in \Z$,
  \[ \operatorname{ord}_{m}\bigl(z \mapsto \Gamma_{\R}(z+a)\bigr) \geq 0
     \iff m+a \text{ is odd or positive}, \]
  since the poles of $\Gamma_{\R}$ are exactly the non-positive even integers.
\end{lemma}
\begin{proof}
  \leanok
  $1/\Gamma_{\R}$ is entire (\lean{Complex.differentiable\_Gammaℝ\_inv}), so
  $z \mapsto \Gamma_{\R}(z+a)$ is the inverse of an entire function $g$ --- pointwise, including
  at the poles, where both sides are $0$ because $0^{-1} = 0$ --- and its order at $s$ is
  $-\operatorname{ord}_s g \leq 0$. That order is finite: $g$ is entire and $g(1-a) = 1$, so by
  the identity theorem $g$ does not vanish identically near $s$. Hence the order is $-n$ for a
  natural number $n$, and $n \neq 0$ if and only if $g(s) = 0$, i.e. if and only if
  $\Gamma_{\R}(s+a) = 0$; at an integer that is \lean{Complex.Gammaℝ\_intCast\_eq\_zero\_iff}.
\end{proof}

\begin{lemma}[A product of $\Gamma_{\R}$'s is regular iff every factor is]
  \label{lem:prodGammaR}
  \lean{Complex.prodGammaℝ, Complex.meromorphicOrderAt\_prodGammaℝ\_nonpos,
    Complex.meromorphicOrderAt\_prodGammaℝ\_nonneg\_iff,
    Complex.meromorphicOrderAt\_prodGammaℝ\_intCast\_nonneg\_iff}
  \leanok
  \uses{lem:GammaR-pole}
  Write $\gamma_{\mathrm{shifts}}(s) = \prod_{a \in \mathrm{shifts}} \Gamma_{\R}(s+a)$. Its
  order at any point is at most $0$, and
  \[ \operatorname{ord}_{s}\gamma_{\mathrm{shifts}} \geq 0
     \iff \operatorname{ord}_{s}\Gamma_{\R}(\cdot + a) \geq 0 \text{ for every }
       a \in \mathrm{shifts}. \]
  At an integer $m$ this says that $\gamma_{\mathrm{shifts}}$ is regular exactly when $m+a$ is odd
  or positive for every shift $a$.
\end{lemma}
\begin{proof}
  \leanok
  Induction on the multiset. The order of a product is the sum of the orders; each factor has
  order $\leq 0$ by Lemma~\ref{lem:GammaR-pole}, and a sum of non-positive elements of
  $\Z \cup \{\infty\}$ is $\geq 0$ exactly when both summands are. The base case is the constant
  $1$, of order $0$.
\end{proof}

\begin{remark}[Criticality is computed, not chosen]
  \label{rem:isCritical-derived}
  \uses{def:isCritical, lem:prodGammaR, rem:crit-spec}
  Deligne defines $n$ to be critical for $M$ when neither $L_{\infty}(M,s)$ nor
  $L_{\infty}(M^{\vee},w+1-s)$ has a pole at $n$ --- criticality is not independent data, it is a
  consequence of the archimedean factor, which is itself determined by the Hodge numbers and
  $F_{\infty}$. Definition~\ref{def:isCritical} says exactly that, with one simplification:
  it writes $\gamma(M,w+1-n)$ rather than $\gamma(M^{\vee},w+1-n)$, which is legitimate because
  the package's $L$ is Deligne's and the two archimedean factors agree there --- for GL(1) by
  Lemma~\ref{lem:gammaFactor-inv}.

  Two successive versions of Definition~\ref{def:pkg} got this wrong in ways worth recording,
  because each fix shrinks what a reader has to take on faith.

  The first carried $\mathrm{IsCritical}$ as a \emph{field}: an arbitrary predicate the instance
  supplied, constrained by nothing. A reader could only check it by reading the instance's
  characterisation theorem and agreeing with it.

  The second supplied instead a bare function $\gamma \colon \Omega \to \C \to \C$ and read ``no
  pole at $n$'' as ``$\gamma(M,n) \neq 0$'', on the grounds that mathlib's \lean{Complex.Gamma}
  takes the value $0$ at its poles. That went wrong twice over. The reading is \emph{false} for a
  general function --- a zero need not be a pole, and an instance supplying $\gamma(M,s) = s-3$
  would have the framework declare $3$ non-critical on account of an honest zero --- so the
  critical set was still chosen by the instance, merely through a zero set rather than through a
  predicate. And even where the reading is sound it is a paraphrase: a definition that means
  ``no pole'' should say ``no pole''.

  Definition~\ref{def:pkg} now supplies the Hodge data itself: a multiset of integer shifts,
  from which $\gamma$ is \emph{defined} as $\prod_a \Gamma_{\R}(s+a)$. Every archimedean
  $L$-factor has this shape --- Deligne's recipe writes $L_{\infty}(M,s)$ as a product of
  $\Gamma_{\C}(s-p)$ over the $h^{p,q}$ with $p<q$ and of $\Gamma_{\R}(s-p)$,
  $\Gamma_{\R}(s-p+1)$ over the $F_{\infty}$-eigenspaces on the $p=q$ part, and
  $\Gamma_{\C}(s) = \Gamma_{\R}(s)\Gamma_{\R}(s+1)$
  (\lean{Complex.Gammaℝ\_mul\_Gammaℝ\_add\_one}) splits each $\Gamma_{\C}$ in two. Criticality
  is then stated as what it is: $\operatorname{ord}_{n}\gamma \geq 0$, in mathlib's
  \lean{meromorphicOrderAt}, with no appeal to a junk-value convention at all. That the
  condition is nevertheless \emph{checkable} is Lemma~\ref{lem:prodGammaR}, which turns it into
  arithmetic on the shifts.
  The residual trust is exactly the Hodge numbers, and that is what Remark~\ref{rem:crit-spec} is
  for: for GL(1) it is Lemma~\ref{lem:gl1-crit-pkg}, checked against mathlib's own
  \lean{DirichletCharacter.gammaFactor} by Lemma~\ref{lem:gl1-gammaFactor-pkg}.
\end{remark}

\begin{remark}[Why the weight is data, and what reflecting about $1-n$ would cost]
  \label{rem:weight}
  \uses{def:pkg, def:isCritical}
  Deligne's functional equation for a motive of weight $w$ exchanges $s$ and $w+1-s$, and
  Definition~\ref{def:isCritical} reflects about that point. An earlier version hard-coded
  $1-n$, i.e. $w = 0$. That is correct for GL(1), whose Artin motives all have weight $0$, and it
  is wrong --- silently, and in the direction that makes the conjecture say nothing --- for every
  odd-weight motive.

  Take $M = H^{1}(E)$ for an elliptic curve: $\gamma(s) = \Gamma_{\C}(s)$, with poles at
  $s \in \Z_{\leq 0}$, and $w = 1$. The correct condition --- $\gamma$ regular at $n$ and at
  $2-n$ --- gives $n \geq 1$ and $n \leq 1$, the single critical point $s=1$. Reflecting about
  $1-n$ instead gives $n \geq 1$ and $n \leq 0$: empty. The instance would then land in
  Theorem~\ref{thm:vacuous-conjecture} and Deligne's conjecture would be vacuously true for every
  elliptic curve.

  Nor can this be absorbed into the normalisation of $L$. Recentring the functional equation on
  $1/2$ costs a shift by $w/2$: for even $w$ that is an integer and harmless, but for odd $w$ it
  moves the critical points into $\Z + \tfrac12$, where an integer $n$ cannot see them, and moves
  the Gamma shifts out of $\Z$ with them. For a weight-$k$ modular form, whose motive has weight
  $k-1$, the hard-coded version covers odd $k$ and excludes every even $k$ --- elliptic curves
  among them. Carrying $w$ as data costs GL(1) one field set to $0$ and keeps
  Definition~\ref{def:pkg} usable for GL(2).

  The weight belongs to the package rather than to the object: it is a \emph{common} weight,
  which is what makes $M_1 \oplus M_2$ pure and Definition~\ref{def:sum} meaningful. A collection
  whose members have different weights is split into one package per weight.
\end{remark}

\begin{remark}[Why (x) is part of the definition, and why nondegeneracy is not]
  \label{rem:not-a-pkg}
  \uses{def:pkg, def:hasCriticalInteger}
  Requirement (x) is not a side condition bolted on afterwards. Without it the package with
  $c^{+} \coloneqq 0$ has $\widetilde{L}_F \equiv 0$, since $x/0 = 0$ in Lean, so both halves
  of Definition~\ref{def:conjecture} hold for free --- $0$ lies in every subfield of $\C$ and
  is fixed by every $\sigma$. Deligne's $c^{+}$ is a determinant of a period matrix and is
  nonzero by construction, so (x) costs a genuine instance nothing, and making it part of the
  definition means the degenerate case cannot be written down at all.

  An earlier version of this definition had a second requirement of the same apparent kind:
  that every object have a critical integer, ruling out $\mathrm{IsCritical} \coloneqq \bot$,
  which likewise satisfies Definition~\ref{def:conjecture} vacuously. That requirement has been
  removed, and the reason is worth recording because it is not a matter of taste. It conflates
  two situations that the shape of the critical set cannot distinguish: an instantiation that
  \emph{dodges} the question by declaring nothing critical, and an object whose
  $\mathrm{IsCritical}$ is \emph{correct} and happens to be empty. The second is not
  pathological --- $h^{0}(\operatorname{Spec} K)$ for $K$ not totally real is such an object,
  and Deligne's conjecture is meaningfully and vacuously true for it
  (Theorem~\ref{thm:vacuous-conjecture}) --- and as a clause of Definition~\ref{def:pkg} the
  requirement made those motives inexpressible. What actually separates the two cases is whether
  $\mathrm{IsCritical}$ is the \emph{right} predicate, and that is already the job of the
  characterisation obligation of Remark~\ref{rem:crit-spec}, which is strictly stronger: it also
  catches a critical set that is wrong and nonempty, which a nonemptiness clause sails past.
  Nondegeneracy is therefore now Definition~\ref{def:hasCriticalInteger}, a predicate discharged
  per instance. It was direct sums (\S\ref{sec:direct-sums}) that forced the issue: the
  requirement is not preserved by $\oplus$.
\end{remark}

\begin{remark}[Why the value field is bounded]
  \label{rem:valueField-bounded}
  \uses{def:pkg, def:algebraicity}
  The value field raises the same question as requirement (x), for the other half of the data.
  Definition~\ref{def:algebraicity} asserts a \emph{membership},
  $\widetilde{L}_{F}(M,n) \in E(M)$, and $E$ is supplied by the instance. Taking
  $E \coloneqq \C$ would make that membership hold for every $L$ and every $c^{+}$, so the
  algebraicity half of Deligne's conjecture would assert nothing at all --- and, unlike a
  vanishing period, a too-large value field is not detectable from the shape of the data: $\C$
  is as much a subfield of $\C$ as $\Q$ is.

  Requirement (ix) is the bound: $E(M)$ consists of algebraic numbers, so the largest field an
  instance may name is $\Qbar$ and not $\C$. There Definition~\ref{def:algebraicity} reads
  ``the normalised critical values are algebraic''. That is a weak statement, but it is a
  statement, which is Theorem~\ref{thm:arithmetic-algebraic}.

  It is weaker than the truth. Deligne's $E(M)$ is a \emph{number field}, finite over $\Q$, and
  requiring finiteness would be the sharp form of the bound; Definition~\ref{def:sum} would
  survive it, a compositum of number fields being a number field. It is not imposed because
  neither instance's statement uses it, and each would then owe a finiteness proof --- for
  GL(1), that $\Q(\chi)$ is finite over $\Q$ --- that buys nothing. What the bound exists for
  is already bought: $E$ can no longer be chosen to make the conjecture true.

  Like requirement (x) and unlike nondegeneracy, the bound passes both tests of
  Remark~\ref{rem:not-a-pkg}. It survives $\oplus$, a compositum of algebraic subfields
  being one; and it excludes no legitimate object, since every motive has a number field for its
  coefficients.
\end{remark}

\begin{theorem}[Algebraicity always asserts algebraicity]
  \label{thm:arithmetic-algebraic}
  \lean{Deligne.Pkg.IsArithmetic.isAlgebraic}
  \leanok
  \uses{def:algebraicity, def:pkg}
  If $F$ satisfies algebraicity at $M$ then $\widetilde{L}_{F}(M,n)$ is algebraic over $\Q$ at
  every critical $n$, whatever $E$ the package supplies.
\end{theorem}
\begin{proof}
  \leanok
  $\widetilde{L}_{F}(M,n) \in E(M)$, and $E(M)$ is algebraic by requirement (ix) of
  Definition~\ref{def:pkg}.
\end{proof}

Two remarks on the shape of this definition, both of which are forced by how the
instances actually want to use it.

\begin{remark}
  $L$ takes a complex argument rather than an integer one. An instance plugs in its
  genuine analytic $L$-function, and the restriction to integer points happens here,
  once. This matters: the proof in every case we know of runs through the functional
  equation, which is a statement about $L$ on all of $\C$.
\end{remark}

\begin{remark}
  The Galois action is a field of the structure rather than a \lean{MulAction}
  instance. An instance almost always wants to \emph{define} its action explicitly
  (for GL(1), $\chi \mapsto \sigma \circ \chi$) rather than to have it found by
  typeclass synthesis, and the action is not canonical in $\Omega$. It is nonetheless a genuine
  action, carried as a monoid homomorphism
  $\operatorname{Gal}(\C/\Q) \to \operatorname{Sym}(\Omega)$ rather than as a bare function together
  with two laws: the compressed form is what mathlib prefers, and it makes $\gal(\sigma)$
  invertible by construction, with inverse $\gal(\sigma^{-1})$, rather than by a derivation.
\end{remark}

\section{The assertions}
\label{sec:assertions}

\begin{definition}[Normalised critical value]
  \label{def:normalizedValue}
  \lean{Deligne.Pkg.normalizedValue}
  \leanok
  \uses{def:pkg}
  For a Deligne package $F$, an object $M$ and $n \in \Z$, set
  \[ \widetilde{L}_F(M,n) \;=\; \frac{L(M,n)}{c^{+}(M(n))}. \]
\end{definition}

\begin{definition}[Algebraicity]
  \label{def:algebraicity}
  \lean{Deligne.Pkg.IsArithmetic}
  \leanok
  \uses{def:normalizedValue}
  $F$ satisfies \emph{algebraicity} at $M$ if for every critical $n$,
  $\widetilde{L}_F(M,n) \in E(M)$.
\end{definition}

\begin{definition}[Galois equivariance]
  \label{def:equivariance}
  \lean{Deligne.Pkg.IsEquivariant}
  \leanok
  \uses{def:normalizedValue}
  $F$ satisfies \emph{equivariance} at $M$ if for every $\sigma \in \operatorname{Gal}(\C/\Q)$
  and every critical $n$,
  \[ \sigma\!\left(\widetilde{L}_F(M,n)\right) \;=\; \widetilde{L}_F(M^{\sigma},n). \]
\end{definition}

Equivariance is the sharp form. It is strictly more information than algebraicity,
and it is the half that pins down the period normalisation: a period that is off by a
transcendental factor can still satisfy a naive algebraicity statement, but it cannot
satisfy equivariance.

\begin{definition}[Deligne's conjecture for a package]
  \label{def:conjecture}
  \lean{Deligne.Pkg.Conjecture}
  \leanok
  \uses{def:algebraicity, def:equivariance}
  $F$ satisfies \emph{Deligne's conjecture} at $M$ if it satisfies both algebraicity and
  equivariance at $M$.
\end{definition}

\begin{lemma}[Rational values are automatically equivariant]
  \label{lem:rational-equivariant}
  \lean{Deligne.Pkg.IsArithmetic.isEquivariant\_of\_valueField\_eq\_bot,
    IntermediateField.apply\_eq\_self\_of\_mem\_bot}
  \leanok
  \uses{def:algebraicity, def:equivariance}
  Let $M$ be an object with $E(M) = \Q$ and $M^{\sigma} = M$ for every
  $\sigma \in \operatorname{Gal}(\C/\Q)$. Then algebraicity at $M$ implies equivariance at $M$.
\end{lemma}
\begin{proof}
  \leanok
  A $\Q$-algebra endomorphism of $\C$ fixes the bottom intermediate field $\Q$ pointwise, so
  $\sigma$ fixes $\widetilde{L}_{F}(M,n)$, which lies in $\Q$ by algebraicity; and
  $M^{\sigma}=M$ makes the right-hand side of equivariance the same number.
\end{proof}

This is not a weakening of Section~\ref{sec:assertions}'s claim that equivariance is strictly
more information than algebraicity. It says only that the gap between them closes when the
value field is $\Q$ \emph{and} the object is Galois-fixed, which is exactly the situation
of a motive over $\Q$ with rational coefficients; for an object with a larger $E(M)$, or one the
action moves, equivariance still carries the extra content.

\section{The period and the \texorpdfstring{$L$}{L}-function}
\label{sec:no-L-field}

\begin{theorem}[The period may not be the $L$-value]
  \label{thm:period-eq-vacuous}
  \lean{Deligne.Pkg.conjecture\_of\_period\_eq\_mul\_L, Deligne.Pkg.conjecture\_of\_L\_eq\_zero}
  \leanok
  \uses{def:conjecture, lem:isCritical-galois}
  Let $F$ be a Deligne package, $M$ an object and $e \colon \Z \to E(M)$.
  If $c^{+}(M^{\sigma}(n)) = \sigma(e(n))\,L(M^{\sigma},n)$ for every
  $\sigma \in \operatorname{Gal}(\C/\Q)$ and every critical $n$ --- that is, if on the Galois
  orbit of $M$ the period is a rescaling of the $L$-value by an element of the value field ---
  then $F$ satisfies Deligne's conjecture at $M$. Likewise if $L(M^{\sigma},n) = 0$ for every
  $\sigma$ and every critical $n$.
\end{theorem}
\begin{proof}
  \leanok
  At a critical $n$ the period does not vanish (requirement (x) of Definition~\ref{def:pkg}),
  so neither does $L(M^{\sigma},n)$, and $\widetilde{L}_{F}(M^{\sigma},n) = \sigma(e(n))^{-1}$;
  Lemma~\ref{lem:isCritical-galois} supplies criticality at $M^{\sigma}$. At $\sigma = 1$ this
  lies in $E(M)$, and $\sigma$ carries $e(n)^{-1}$ to $\sigma(e(n))^{-1}$, which is equivariance.
  For a vanishing $L$-function the normalised value is $0$ on the orbit, which lies in every
  field and is fixed by every $\sigma$.
\end{proof}

Theorem~\ref{thm:period-eq-vacuous} is the one degenerate shape that no requirement of
Definition~\ref{def:pkg} can exclude. An instance defines all its fields at once, so $c^{+}$ may
be written in terms of $L$, and no requirement stated \emph{within} the structure can detect
that it was. It is worth being precise about what the requirements do and do not do.
Requirement (x) forbids a period that \emph{vanishes}; it says nothing about where a nonvanishing
period came from. Uniqueness would not help either: a predicate ``$x$ is the period of $M(n)$''
defined as $x = L(M,n)$ is trivially satisfied by exactly one $x$.

An earlier version of Definition~\ref{def:pkg} removed $L$ from the data for this reason, so
that an instance asserted its conjecture in the form
$\forall L,\ (\text{$L$ is the $L$-function of $M$}) \implies F.\mathrm{Conjecture}(L,M)$,
with the period fixed before $L$ was quantified. That made the independence of $c^{+}$ from $L$
a fact about the types. It has been given up in favour of carrying $L$ as a field, because the
package should be the complete data of the conjecture and because nothing formal prevents an
author from writing $F$ and $L$ in the same file with the same number in mind in either shape.
What remains is a reading obligation: a reader of an instance checks that its period does not
mention its $L$-function. For GL(1) the period is a Gauss sum times a power of $2\pi i$
(Definition~\ref{def:gl1-period}) and the $L$-function is mathlib's; for
Chapter~\ref{chap:sympow} the period is a function of a period lattice and of nothing else
(Definition~\ref{def:periodpair-periods}), and the $L$-function is quantified over outside the
package (Definition~\ref{def:sympow-conjecture}).

\section{Guarding against vacuous instances}
\label{sec:guards}

The framework of \S\ref{def:pkg} buys reuse at a price that must be paid
explicitly. Requirements (ix) and (x) of Definition~\ref{def:pkg} rule out a vanishing period
and a value field too big to say anything, but nothing in the definition rules out an
instance whose $\mathrm{IsCritical}$ is
honest-looking and still wrong. A reader cannot tell, from the fact that \lean{Conjecture}
has been proved, whether the critical set it was proved for is the right one.

We therefore require every instance in this development to discharge two further
obligations alongside its main theorem. These are not part of the statement of
Deligne's conjecture; they are what makes an instantiation of it trustworthy.

\begin{definition}[Nondegeneracy]
  \label{def:hasCriticalInteger}
  \lean{Deligne.Pkg.HasCriticalInteger}
  \leanok
  \uses{def:pkg}
  An object $M$ is \emph{nondegenerate} if it has at least one critical integer.
\end{definition}

\begin{theorem}[The degenerate case is vacuously true]
  \label{thm:vacuous-conjecture}
  \lean{Deligne.Pkg.Conjecture.of\_not\_hasCriticalInteger,
    Deligne.Pkg.conjecture\_of\_forall\_not\_isCritical}
  \leanok
  \uses{def:hasCriticalInteger, def:conjecture}
  If $M$ has no critical integer then $F$ satisfies Deligne's conjecture at $M$.
\end{theorem}
\begin{proof}
  \leanok
  Both halves of Definition~\ref{def:conjecture} quantify over critical integers, so both
  are vacuously true. This is the content of Remark~\ref{rem:not-a-pkg}'s second
  paragraph: the statement is not a loophole but a theorem, and the objects it applies to are
  real. $\zeta_K$ for $K$ with a complex place is one, and its emptiness is itself a theorem
  --- the archimedean factor acquires a $\Gamma_{\C}$, whose poles at $s \leq 0$ and whose
  reflected poles at $s \geq 1$ between them exclude every integer.
\end{proof}

\begin{lemma}[Criticality is Galois-stable]
  \label{lem:isCritical-galois}
  \lean{Deligne.Pkg.isCritical\_galoisAction\_iff}
  \leanok
  \uses{def:pkg, def:isCritical}
  For every $\sigma \in \operatorname{Gal}(\C/\Q)$,
  $\mathrm{IsCritical}(M^{\sigma},n) \iff \mathrm{IsCritical}(M,n)$.
\end{lemma}
\begin{proof}
  \leanok
  Criticality depends on $M$ only through $\mathrm{shifts}(M)$ (Definition~\ref{def:isCritical}),
  and requirement (vii) says $\mathrm{shifts}(M^{\sigma}) = \mathrm{shifts}(M)$.
\end{proof}

\emph{Coherence} --- this lemma together with the transport $E(M^{\sigma}) = \sigma(E(M))$ of the
value field --- was once an obligation each instance discharged separately. Requirements (vii)
and (viii) of Definition~\ref{def:pkg} absorb it: what an instance now proves, in proving those
two requirements, is exactly what it used to prove afterwards.

\begin{remark}[The second obligation]
  \label{rem:crit-spec}
  Nondegeneracy is a property stated within the framework. The second obligation
  cannot be: each instance must prove a \emph{characterisation}
  theorem identifying its $\mathrm{IsCritical}$ predicate with an independently and
  concretely specified condition, stated without reference to the framework. For
  GL(1) this is Lemma~\ref{lem:gl1-crit-spec}, and it is checked against
  $\zeta$ in Remark~\ref{rem:zeta-sanity}. An instance that proves
  \lean{Conjecture} but omits its characterisation theorem has
  not been shown to be about anything.
\end{remark}

\section{Direct sums}
\label{sec:direct-sums}

A framework that could not express the direct sum of two motives would be of no use beyond
its first instance: every Artin motive of rank $>1$ over $\Q$ is a sum of smaller ones, and
$\zeta_K$ is the standard example. This section carries out that construction. It is also what
diagnosed the shape of Definition~\ref{def:pkg}: in an earlier version, nondegeneracy was a
clause of that definition and the construction below could not be carried out without an extra
hypothesis. Nondegeneracy is now Definition~\ref{def:hasCriticalInteger}, the sum is total, and
what remains to record is one genuine limitation that $\oplus$ exposes and the design change does
not remove.

\begin{definition}[Direct sum of packages]
  \label{def:sum}
  \lean{Deligne.Pkg.sum}
  \leanok
  \uses{def:pkg}
  Let $F_1$ on $\Omega_1$ and $F_2$ on $\Omega_2$ be Deligne packages \emph{of the same weight},
  $w_1 = w_2$. Then $F_1 \oplus F_2$ is the package on $\Omega_1 \times \Omega_2$ of weight $w_1$
  with
  \[ L = L_1 L_2, \qquad \mathrm{shifts} = \mathrm{shifts}_1 + \mathrm{shifts}_2, \qquad
     E = E_1 E_2, \qquad c^{+} = c^{+}_1 c^{+}_2, \]
  and the componentwise Galois action. The multiset sum of the shifts is the product
  $\gamma = \gamma_1 \gamma_2$ of the archimedean factors. Equality of the weights is not a
  technicality: the functional equation of $\Lambda_1 \Lambda_2$ exchanges $s$ and $w+1-s$ only
  when the two factors reflect about the same point, which is the statement that
  $M_1 \oplus M_2$ is pure.
\end{definition}
\begin{proof}
  \leanok
  Each clause is the motivic fact: $L$-functions and Deligne periods are multiplicative on
  direct sums (the period because the comparison isomorphism is block diagonal, so its
  determinant is the product); the value field of a sum is the compositum; and the
  archimedean factor of a sum is the product of the factors. Of the requirements of
  Definition~\ref{def:pkg}, (vii) and (viii) hold componentwise, (viii) because $\sigma_{*}$ on subfields
  is a left adjoint (to $\sigma^{*}$) and therefore preserves joins, which is mathlib's
  \lean{IntermediateField.map\_sup}; (ix) because a compositum of algebraic intermediate fields
  lies in $\Qbar$;
  and (x) because a product of nonzero periods is nonzero.

  Criticality is no longer among the choices: it is computed from $\gamma$, and since a product is
  nonzero exactly when both factors are, $n$ is critical for $M_1 \oplus M_2$ if and only if it is
  critical for each summand. What used to be a definitional decision is
  Lemma~\ref{lem:isCritical-sum}.
\end{proof}

\begin{remark}[Nondegeneracy is not inherited, and that is now harmless]
  \label{rem:sum-obstruction}
  \uses{def:sum, def:hasCriticalInteger}
  Definition~\ref{def:hasCriticalInteger} is \emph{not} preserved by $\oplus$: two motives may
  each have critical integers and yet share none, so $M_1$ and $M_2$ may both be nondegenerate
  while $M_1 \oplus M_2$ is not. In the GL(1) case this is exactly a parity condition
  (Lemma~\ref{lem:gl1-common-critical}), and it fails for the pair $(\mathbf{1}, \chi)$ with
  $\chi$ odd. That pair is not a curiosity: its sum is $h^{0}(\operatorname{Spec} K)$ for
  $K$ imaginary quadratic, whose $L$-function is
  $\zeta_K = \zeta \cdot L(\chi_{-d}, \cdot)$ --- and $\zeta_K$ genuinely has no critical
  integers, because the archimedean factor acquires a $\Gamma_{\C}$, whose poles at $s \leq 0$
  and whose reflected poles at $s \geq 1$ between them exclude every integer.

  With nondegeneracy a predicate rather than a clause, none of this obstructs anything.
  $F_1 \oplus F_2$ is a Deligne package regardless; $h^{0}(\operatorname{Spec} K)$ is an object
  of it; and Deligne's conjecture holds there by Theorem~\ref{thm:vacuous-conjecture} rather than
  being unstatable. What is lost is only that an instance built from a sum must prove
  nondegeneracy for itself when it wants it, which is the correct state of affairs: for
  $h^{0}(\operatorname{Spec} K)$ the honest theorem is that nondegeneracy \emph{fails} unless
  $K$ is totally real.

  A second limitation surfaces in the same construction and is worth separating from the first.
  Definition~\ref{def:pkg} makes the value field a \emph{field}
  (\lean{IntermediateField} $\Q\ \C$), but the value object of a direct sum is properly the
  \emph{ring} $E_1 \times E_2$: Deligne's conjecture for $M_1 \oplus M_2$ is really the pair of
  conjectures, one over each factor. Definition~\ref{def:sum} uses the compositum $E_1 E_2$
  instead, which is correct but lossy --- $L_1L_2/(c_1^{+}c_2^{+}) \in E_1E_2$ is implied by,
  and strictly weaker than, $L_1/c_1^{+} \in E_1$ together with $L_2/c_2^{+} \in E_2$. So
  Theorem~\ref{thm:conjecture-sum} concludes something weaker than what it consumes. Nothing is
  wrong with it; but a framework that intends to reach beyond a single instance should know that
  the field is a constraint. Unlike the nondegeneracy clause, this one has not been removed:
  restating the conjecture over a value \emph{ring} would touch every definition in
  \S\ref{sec:assertions}, and every instance in sight has an honest field.
\end{remark}

\begin{lemma}[Criticality for a direct sum]
  \label{lem:isCritical-sum}
  \lean{Deligne.Pkg.isCritical\_sum\_iff, Deligne.Pkg.gammaFactor\_sum}
  \leanok
  \uses{def:sum, def:isCritical}
  $n$ is critical for $M_1 \oplus M_2$ if and only if it is critical for $M_1$ and for $M_2$.
\end{lemma}
\begin{proof}
  \leanok
  The shifts of the sum are the multiset union of the shifts, so by
  Lemma~\ref{lem:prodGammaR} the sum is regular at a point exactly when both summands are; apply
  this at $n$ and at $w+1-n$, which is the same point for both summands because the weights
  agree.
\end{proof}

\begin{lemma}[Normalised values multiply]
  \label{lem:normalizedValue-sum}
  \lean{Deligne.Pkg.normalizedValue\_sum, Deligne.Pkg.galoisAction\_sum}
  \leanok
  \uses{def:sum, def:normalizedValue}
  The normalised value of $F_1 \oplus F_2$ at $(M_1, M_2)$ and $n$ is the product of the
  normalised values of $M_1$ and $M_2$ at $n$.
\end{lemma}
\begin{proof}
  \leanok
  Both the $L$-value and the period are products, and $(ac)/(bd) = (a/b)(c/d)$.
\end{proof}

\begin{theorem}[Deligne's conjecture is inherited by direct sums]
  \label{thm:conjecture-sum}
  \lean{Deligne.Pkg.Conjecture.sum, Deligne.Pkg.IsArithmetic.sum,
    Deligne.Pkg.IsEquivariant.sum}
  \leanok
  \uses{def:sum, lem:normalizedValue-sum, def:conjecture, def:algebraicity, def:equivariance}
  If $F_1$ satisfies the conjecture at $M_1$ and $F_2$ at $M_2$, then $F_1 \oplus F_2$
  satisfies it at $(M_1, M_2)$.
\end{theorem}
\begin{proof}
  \leanok
  \uses{lem:normalizedValue-sum}
  A critical integer for the sum is critical for each summand, so
  Lemma~\ref{lem:normalizedValue-sum} writes the normalised value as a product of two values
  known to lie in $E_1$ and $E_2$ respectively; both lie in the compositum, and so does the
  product. For equivariance, $\sigma$ is multiplicative and each factor transports by
  hypothesis. The one Lean wrinkle is that the transported object
  $\sigma \cdot (M_1,M_2)$ is only \emph{definitionally} the pair
  $(\sigma \cdot M_1, \sigma \cdot M_2)$, so Lemma~\ref{lem:normalizedValue-sum} has to be
  instantiated at it by hand rather than rewritten with.
\end{proof}

\section{Galois descent}
\label{sec:descent}

Section~\ref{sec:direct-sums} produces, for a sum of two motives, the weakest true statement:
$\widetilde{L} \in E_1E_2$. The statement one actually wants for $\zeta_K$ is $\widetilde{L} \in
\Q$, which is strictly stronger --- $\Q$ is the smaller field --- and it is not a formal
consequence of the two summand statements. It follows instead from \emph{equivariance}, which up
to this point in the development has been proved and then used for nothing.

\begin{theorem}[Descent]
  \label{thm:descent}
  \lean{Deligne.Pkg.map\_prod\_normalizedValue}
  \leanok
  \uses{def:equivariance, def:normalizedValue}
  Let $(M_i)_{i \in I}$ be a finite family of objects, $n$ critical for each $M_i$, and suppose
  each $M_i$ satisfies equivariance. If $\sigma \in \operatorname{Gal}(\C/\Q)$ permutes the family
  --- that is, there is a permutation $e$ of $I$ with $M_i^{\sigma} = M_{e(i)}$ --- then
  \[ \sigma\Bigl(\prod_{i \in I} \widetilde{L}_F(M_i, n)\Bigr)
     = \prod_{i \in I} \widetilde{L}_F(M_i, n). \]
\end{theorem}
\begin{proof}
  \leanok
  Apply $\sigma$ termwise. Equivariance turns $\sigma(\widetilde{L}_F(M_i,n))$ into
  $\widetilde{L}_F(M_i^{\sigma},n) = \widetilde{L}_F(M_{e(i)},n)$, and reindexing along the
  permutation $e$ returns the original product.
\end{proof}

\begin{lemma}[The product lies in the compositum]
  \label{lem:prod-mem-compositum}
  \lean{Deligne.Pkg.prod\_normalizedValue\_mem\_iSup}
  \leanok
  \uses{def:algebraicity}
  Under the same hypotheses with algebraicity in place of equivariance,
  $\prod_i \widetilde{L}_F(M_i,n)$ lies in the compositum $\bigvee_i E(M_i)$.
\end{lemma}
\begin{proof}
  \leanok
  Each factor lies in its own value field, hence in the join, and a subfield is closed
  under products.
\end{proof}

\begin{remark}[Where this stops, and why]
  \label{rem:descent-gap}
  \uses{thm:descent}
  The fixed field of $\operatorname{Gal}(\C/\Q)$ is $\Q$: an automorphism fixes $\Q$ pointwise,
  and every irrational --- algebraic or transcendental --- is moved by some automorphism. So
  Theorem~\ref{thm:descent} says the product is \emph{rational}, and that is the Klingen--Siegel
  shape: for $K$ abelian and totally real, $\zeta_K = \prod_{\chi \in X} L(\chi, \cdot)$ over a
  Galois-stable set $X$ of even characters, each factor lies in $\Q(\chi)$, and the product lies
  in $\Q$ --- which does not follow from the factors alone.

  That last step is not formalised, and the two things missing are both absent from mathlib, not
  from this development. First, mathlib has no statement that the fixed field of
  $\operatorname{Gal}(\C/\Q)$ is $\Q$, so Theorem~\ref{thm:descent} stops at the invariance equation.
  Second, mathlib has no \lean{Fintype} instance for Dirichlet characters, so GL(1) has no
  \emph{nontrivial} Galois-stable family to instantiate the theorem at. A trivial one does exist,
  and it is worth saying so: any Galois-fixed object gives a Galois-stable singleton, and the
  level-one character is fixed (\lean{MulChar.ringHomComp\_one}), so the hypotheses are
  satisfiable. What is unavailable is a family about which the theorem says something new. Both
  gaps are stated precisely so a later pass can check whether they have closed rather than
  rediscover the obstruction.

  The hypothesis of Theorem~\ref{thm:descent} still hands over a permutation of the index rather
  than deriving one, and that is now a convenience rather than a necessity: since $\gal$ is a
  group action, $\gal(\sigma)$ is a bijection with inverse $\gal(\sigma^{-1})$, so a family
  stable under it and indexed injectively determines the permutation. Deriving it needs the
  injectivity of the indexing together with a counting argument, and is left until a caller wants
  the weaker hypothesis. The value field of Remark~\ref{rem:sum-obstruction} remains the one
  limitation of Definition~\ref{def:pkg} that a second instance should be expected to press
  on.

  Two consistency checks, from a direction independent of the $\Gamma$-factor computation. For $K$
  abelian and totally real every $\chi \in X$ is even, so by Lemma~\ref{lem:gl1-common-critical}
  the critical sets coincide and the product is nondegenerate. For $K$ with a complex place $X$
  contains an odd character, the critical sets are disjoint, and there is no critical $n$ at all
  --- matching Remark~\ref{rem:sum-obstruction}.
\end{remark}
